% ============================================================================
% SEC05.TEX - Section 7.5: Blend Trees
% ============================================================================

\section{Blend Trees}

\begin{learningoutcomes}
\item Memahami konsep Blend Tree
\item Membuat Blend Tree untuk smooth transitions
\item Menggunakan 1D dan 2D Blend Trees
\end{learningoutcomes}

\subsection{Pengertian Blend Tree}

\begin{definisi}[Blend Tree]
Blend Tree adalah struktur dalam Animator Controller yang memungkinkan blending halus antar animasi berdasarkan parameter \cite{UnityAnimation2024}.
\end{definisi}

\subsection{1D Blend Tree}

Untuk blending berdasarkan satu parameter (misalnya: speed):
\begin{itemize}
\item Idle (speed = 0)
\item Walk (speed = 1)
\item Run (speed = 2)
\end{itemize}

\subsection{2D Blend Tree}

Untuk blending berdasarkan dua parameter (misalnya: horizontal dan vertical movement):
\begin{itemize}
\item Idle (0, 0)
\item Forward (0, 1)
\item Backward (0, -1)
\item Left (-1, 0)
\item Right (1, 0)
\item Forward-Left (-0.7, 0.7)
\item dll
\end{itemize}

\begin{catatan}
Blend Tree memberikan transisi yang lebih smooth antara animasi. Gunakan untuk movement animations yang perlu blending natural.
\end{catatan}

\subsection{Ringkasan Section}

\begin{itemize}
\item Blend Tree memungkinkan blending halus antar animasi
\item 1D Blend Tree untuk satu parameter
\item 2D Blend Tree untuk dua parameter
\item Berguna untuk movement animations yang perlu transisi natural
\end{itemize}
